\section{Conclusion}
\label{sec:conclusion}

A phishing verdict is acted on, contested, and audited, yet the accuracy that has
become the field's headline measure speaks to none of those moments. This paper
takes one route toward narrowing that gap: read a detector's reliability from the
\emph{evidence} a verdict rests on rather than from the label it emits. We
instantiate the idea as \toolname{}, which trusts a verdict in proportion to how
much a panel of independent agents agree on cues the page itself confirms, returns
those cues as the explanation, and abstains otherwise. Three findings carry the
contribution beyond the headline: (i) moving the unit of agreement from labels to
grounded evidence is what earns the improvement -- it gives competitive area under the
risk--coverage curve and, by a clear and statistically significant margin, the
best-calibrated trust score and the highest selective accuracy, where
using several models alone does not; (ii) cross-family \emph{diversity}, not the agent
count, is the dominant design variable -- collapsing the panel to a same-family one
degrades both the ranking and full-coverage accuracy sharply; and (iii) calibrating the
grounded score against real labels makes its reported trust directly actionable, the
clearest separation from every baseline. Trust and explanation, in short, can be earned
per verdict and decoupled from the raw accuracy the detector is otherwise held to.

\sstitle{Limitations}
\label{sec:limitations}
The evaluation is scoped and several results are weaker than a clean design would hope.
\emph{(a) Benchmark.} The results rest on a single corpus of feed-sourced phishing pages
versus legitimate login pages ($998$ test) under a fixed three-model panel, so the
headline numbers depend on corpus composition and the specific agents. Two checks bound this
dependence: a leave-brands-out test (\autoref{sec:exp_brands}) shows the calibrated trust
transfers to entirely unseen brands with no loss, and swapping the $3$B text agents for $7$B/$8$B
siblings does \emph{not} improve detection (\autoref{sec:exp_negatives}) -- the bottleneck is
cue elicitation, not model scale -- so the small-model panel is not the limiting factor.
Live cross-feed validation on freshly collected pages nonetheless remains future work. \emph{(b)
The right-label-wrong-reason effect is modest here.} Because this corpus's phishing pages
carry an unambiguous off-brand host and few \emph{shared hallucinations}, the perturbation
fragility that GEA is designed to flag appears only on the harder lookalike subset and
only weakly ($1.2\times$ low-vs-high-GEA); the mechanism is sound but its payoff is
benchmark-dependent. \emph{(c) Grounding adds little to the ranking on \emph{clean} data} for the
same reason -- the agreed cues almost always survive their tool check, and on this corpus
agreement and groundedness are strongly correlated (\autoref{sec:form-gea}), so
$\agr{\cdot}\grd$ ranks much like $\agr$ alone. Groundedness earns its place in the regime the corpus does
not contain: when we synthesize evidence corruption (\autoref{sec:exp_adv}, RQ7) by
cloaking the form action or removing the brand logo, $\grd$ falls and the detector
abstains rather than acts on the corrupted cue, while label-confidence baselines stay
overconfident -- the fail-safe role the clean-corpus AURC cannot reveal. A naturally
occurring benchmark richer in shared hallucinations would exercise it without synthetic
perturbation. We traced this to the benchmark's
\emph{evidence} composition rather than to a design fault, ruling out two cheaper
explanations. Weakening the panel does not help: a deliberately weaker all-local-3B
panel (including a 3B vision model) surfaces no confident-but-ungrounded consensus -- it
merely \emph{lowers} agreement (accuracy near chance), so there is nothing for grounding to
catch. Sharpening the domain signal does not help either: restricting to the
lookalike/typosquat subset (brand token embedded in the host) leaves grounding neutral
(A-only AURC $1.93$ vs.\ $2.05$ for $A{\cdot}G$), because those pages still carry a genuine
brand logo and are detected on that visual evidence ($86\%$ caught) -- the verdict is robust,
not fragile. The limitation is therefore intrinsic to a corpus whose phishing pages always
present a reliable visual brand cue; exercising grounding and the fragility gradient would
require evasive, minimal-cue phishing (cloaked or logo-free pages resting on a single brittle
signal), which this corpus does not contain at any domain composition. \emph{(d) Tooling.} We re-run D1
(Phishpedia) in a containerized \texttt{detectron2} on the full $998$-page test split, and the
heavier D2 (PhishIntention) is cited rather than re-run.
Our main runs ground the brand cue by a page-content check; we do show that D1 can be
substituted as that tool with no loss of ranking (AURC $2.33$ vs.\ $2.29$), so the
``subsumes the detector'' property is measured for D1, but not yet for D2 or the LLM checker. A verification tool is
also sound only for what it checks, so groundedness bounds rather than eliminates gauge
error, and a fully cloaked page carrying no verifiable cue can only be abstained on, not
resolved.

Each bound is also an opening. Strengthening the verifiers toward semantic
plausibility -- not just structural soundness -- would tighten the grounding guarantee
against true-yet-misleading evidence; learning the evidence schema and inferring panel
diversity from opaque third-party models would let the method track new impersonation
tactics without re-engineering the panel; and faster page-internal grounding would
shrink the cloaked-page blind spot rather than route around it. Each direction extends
the evidence-agreement framework without redesigning it, sharpening when a verdict can
be trusted rather than how often it is right.
